% ==========================================
% BẢNG PHÂN CÔNG CÔNG VIỆC
% ==========================================
\section*{BẢNG PHÂN CÔNG CÔNG VIỆC VÀ MỨC ĐỘ ĐÓNG GÓP}
\addcontentsline{toc}{section}{Bảng phân công công việc}

\subsection*{Tóm tắt mã công việc}

Các mã công việc được dùng trong bảng phân công nhằm rút gọn nhiệm vụ của từng thành viên.
Nhóm mã F tương ứng với các nội dung cài đặt và minh họa lý thuyết ở Phần 1; nhóm mã T tương
ứng với các nội dung xử lý dữ liệu thực tế, xây dựng mô hình và hoàn thiện báo cáo ở Phần 2.

\begingroup
\small
\renewcommand{\arraystretch}{1.15}
\begin{table}[H]
\centering
\caption{Tóm tắt các mã công việc Phần 1}
\label{tab:assignment_f_summary}
\begin{adjustbox}{max width=\textwidth}
\begin{tabular}{|p{0.1\textwidth}|p{0.86\textwidth}|}
\hline
\textbf{Mã} & \textbf{Nội dung chính} \\ \hline
F1 & Cài đặt \texttt{ols\_fit}: giải hồi quy tuyến tính OLS bằng Normal Equations, tính hệ số, dự đoán và phần dư. \\ \hline
F2 & Cài đặt \texttt{hat\_matrix}: tính ma trận Hat và kiểm tra các tính chất đối xứng, lũy đẳng, rank, eigenvalues. \\ \hline
F3 & Cài đặt \texttt{model\_metrics}: tính RSS, TSS, MSS, $R^2$, Adj-$R^2$, F-statistic, p-value, MAE và RMSE. \\ \hline
F4 & Cài đặt \texttt{coef\_inference}: tính sai số chuẩn, $t$-statistic, p-value và khoảng tin cậy 95\% cho hệ số. \\ \hline
F5 & Cài đặt \texttt{vif}: tính Variance Inflation Factor để phát hiện đa cộng tuyến giữa các biến độc lập. \\ \hline
F6 & Cài đặt \texttt{ridge\_fit}: Ridge Regression bằng nghiệm dạng đóng, có chuẩn hóa dữ liệu và chọn tham số phạt. \\ \hline
F7 & Cài đặt \texttt{lasso\_fit}: Lasso Regression bằng Coordinate Descent và soft-thresholding. \\ \hline
F8 & Cài đặt \texttt{residual\_plots}: vẽ các biểu đồ chẩn đoán phần dư để kiểm tra giả định của mô hình. \\ \hline
F9 & Cài đặt \texttt{kfold\_cv}: thực hiện k-fold cross-validation để đánh giá mô hình và chọn $\lambda$ tối ưu. \\ \hline
F10 & Mô phỏng Monte Carlo minh họa định lý Gauss-Markov, kiểm tra tính không chệch và phương sai của OLS. \\ \hline
F11 & Tổng hợp notebook Phần 1, minh họa các hàm, so sánh với thư viện và đưa kết quả vào báo cáo. \\ \hline
\end{tabular}
\end{adjustbox}
\end{table}

\begin{table}[H]
\centering
\caption{Tóm tắt các mã công việc Phần 2}
\label{tab:assignment_t_summary}
\begin{adjustbox}{max width=\textwidth}
\begin{tabular}{|p{0.1\textwidth}|p{0.86\textwidth}|}
\hline
\textbf{Mã} & \textbf{Nội dung chính} \\ \hline
T1 & Chọn và tải tập dữ liệu thực tế cho bài toán hồi quy, bảo đảm đủ số quan sát, số biến và có dữ liệu khuyết thiếu. \\ \hline
T2 & Thực hiện EDA: thống kê mô tả, báo cáo missing value, kiểm tra trùng lặp, histogram, boxplot, heatmap và scatter plot. \\ \hline
T3 & Xây dựng \texttt{DataPipeline}: xử lý missing value, ngoại lai, encoding, scaling, polynomial features và tránh data leakage. \\ \hline
T4 & Huấn luyện và so sánh các mô hình OLS đầy đủ, OLS chọn biến, Ridge và Lasso trên tập dữ liệu thực tế. \\ \hline
T5 & Phân tích kết quả mô hình, phân tích phần dư, tổng hợp bảng biểu và hoàn thiện nội dung báo cáo Phần 2. \\ \hline
T6 & Phân tích Feature Importance, trực quan hóa các biến quan trọng và hỗ trợ tổng hợp notebook Phần 2. \\ \hline
T7 & Cài đặt kỹ thuật nâng cao hoặc phần bonus, đồng thời rà soát chất lượng mã nguồn và unit test. \\ \hline
\end{tabular}
\end{adjustbox}
\end{table}
\endgroup

\renewcommand{\arraystretch}{1.5}
\begin{longtable}{|p{0.05\textwidth}|p{0.18\textwidth}|>{\linespread{1.3}\selectfont\arraybackslash}p{0.6\textwidth}|p{0.08\textwidth}|}
    
    \caption{Bảng phân công công việc và mức độ đóng góp} \\
    
    % --- 1. TIÊU ĐỀ Ở TRANG ĐẦU TIÊN ---
    \hline
    \centering \textbf{STT} & \textbf{Họ tên / MSSV} & \textbf{Nhiệm vụ đảm nhiệm chi tiết} & \textbf{Đóng góp} \\ \hline
    \endfirsthead
    
    % --- 2. TIÊU ĐỀ LẶP LẠI Ở CÁC TRANG SAU ---
    \hline
    \centering \textbf{STT} & \textbf{Họ tên / MSSV} & \textbf{Nhiệm vụ đảm nhiệm chi tiết} & \textbf{Đóng góp} \\ \hline
    \endhead
    
    % --- 3. CHÂN BẢNG Ở CÁC TRANG BỊ NGẮT ---
    \hline
    \endfoot
    
    % --- 4. CHÂN BẢNG Ở TRANG CUỐI CÙNG ---
    \hline
    \endlastfoot

    % ================= NỘI DUNG BẢNG =================
    
    % Thành viên 1
    \centering 1 & 
    24120018 \newline \textbf{Đào Thị Quỳnh Anh} & 
    \vspace{-0.4cm}
    \begin{itemize}[itemsep=0.15cm, topsep=0cm, partopsep=0cm]
        \item \textbf{Phần 1:} Cài đặt thuật toán hồi quy có cực tiểu hóa hình phạt: Ridge Regression (nghiệm dạng đóng) và Lasso Regression (giải bằng phương pháp Coordinate Descent).
        \item \textbf{Phần 1:} Xây dựng hệ thống chẩn đoán mô hình trực quan (vẽ 4 biểu đồ phần dư) và cài đặt thuật toán kiểm chứng chéo K-Fold Cross-Validation từ đầu để tinh chỉnh siêu tham số $\lambda$.
        \item \textbf{Phần 2 (T7 \& QA):} Cài đặt mô hình nâng cao (Kernel Ridge Regression hoặc Bayesian Linear Regression) để so sánh đối chiếu. Thực hiện rà soát chất lượng mã nguồn (Quality Assurance), kiểm tra Unit Test và đảm bảo tính nhất quán của hằng số \texttt{RANDOM\_STATE}.
    \end{itemize} \vspace{-0.2cm} & 
    \centering 100\% \tabularnewline \hline
    
    % Thành viên 2
    \centering 2 & 
    24120347 \newline \textbf{Phan Lê Đăng Khoa} & 
    \vspace{-0.4cm}
    \begin{itemize}[itemsep=0.15cm, topsep=0cm, partopsep=0cm]
        \item \textbf{Phần 1:} Lập trình lõi thuật toán Hồi quy tuyến tính đa biến (OLS) bằng phương trình chuẩn (Normal Equations) sử dụng Python thuần.
        \item \textbf{Phần 1:} Tính toán ma trận chiếu (Hat Matrix), kiểm chứng tính chất đại số tuyến tính (lũy đẳng, đối xứng) và trích xuất trị riêng. Cài đặt hệ thống đo lường với các chỉ số thống kê cốt lõi ($R^2$, Adj-$R^2$, RSS, F-statistic).
        \item \textbf{Phần 2 (T3):} Xây dựng lớp \texttt{DataPipeline} hướng đối tượng nhằm tự động hóa các khâu: nội suy KNN, xử lý ngoại lai (Winsorization), chuẩn hóa (Scaling) và mã hóa (Encoding), đảm bảo tuyệt đối không xảy ra rò rỉ dữ liệu (No Data Leakage).
    \end{itemize} \vspace{-0.2cm} & 
    \centering 100\% \tabularnewline \hline
    
    % Thành viên 3
    \centering 3 & 
    24120348 \newline \textbf{Hà Đăng Khôi} & 
    \vspace{-0.4cm}
    \begin{itemize}[itemsep=0.15cm, topsep=0cm, partopsep=0cm]
        \item \textbf{Cơ sở hạ tầng:} Khởi tạo kiến trúc dự án (\texttt{README.md}, \texttt{requirements.txt}) và thiết kế Jupyter Notebook trình bày cơ sở lý thuyết Toán học cho các thuật toán từ F1 đến F5.
        \item \textbf{Phần 2 (T1, T2):} Chọn lọc và tải tập dữ liệu thực tế. Thực hiện Khảo sát dữ liệu toàn diện (EDA): báo cáo dữ liệu khuyết thiếu, phân tích phân phối (Histogram, Boxplot), và tính toán ma trận tương quan (Heatmap).
        \item \textbf{Phần 2 (T6):} Xây dựng thuật toán nhận diện điểm ngoại lai bằng phương pháp IQR. Trực quan hóa mức độ quan trọng của các biến độc lập (Feature Importance) bằng Bar chart và tổng hợp báo cáo Notebook.
    \end{itemize} \vspace{-0.2cm} & 
    \centering 100\% \tabularnewline \hline

    % Thành viên 4
    \centering 4 & 
    24120352 \newline \textbf{Lương Trung Kiên} & 
    \vspace{-0.4cm}
    \begin{itemize}[itemsep=0.15cm, topsep=0cm, partopsep=0cm]
        \item \textbf{Cơ sở hạ tầng:} Thiết lập khung báo cáo \LaTeX{} chuẩn học thuật. Thiết kế Jupyter Notebook minh họa lý thuyết phần chẩn đoán và xác thực chéo (các hàm từ F6 đến F10).
        \item \textbf{Phần 2 (T5):} Lập bảng tổng hợp, so sánh hiệu năng các mô hình trên tập kiểm thử thông qua độ đo MAE, RMSE, $R^2$. Phân tích chuyên sâu biểu đồ phần dư của mô hình tối ưu.
        \item \textbf{Báo cáo:} Đưa ra các nhận định đánh giá giả định Gauss-Markov dựa trên số liệu. Chịu trách nhiệm viết phần Kết luận, căn chỉnh định dạng \LaTeX, xử lý caption hình ảnh và biên soạn tài liệu tham khảo theo chuẩn BibTeX.
    \end{itemize} \vspace{-0.2cm} & 
    \centering 100\% \tabularnewline \hline
    
    % Thành viên 5
    \centering 5 & 
    24120418 \newline \textbf{Nguyễn Minh Quân} & 
    \vspace{-0.4cm}
    \begin{itemize}[itemsep=0.15cm, topsep=0cm, partopsep=0cm]
        \item \textbf{Phần 1:} Cài đặt module suy diễn thống kê: tính toán sai số chuẩn, $t$-statistic, $p$-value và khoảng tin cậy 95\% cho các hệ số hồi quy.
        \item \textbf{Phần 1:} Cài đặt thuật toán tính Hệ số nhân phương sai (VIF) nhằm phát hiện đa cộng tuyến. Lập trình mô phỏng Monte Carlo ($N=1000$) để kiểm chứng thực nghiệm định lý Gauss-Markov.
        \item \textbf{Phần 2 (T4):} Sử dụng Pipeline để huấn luyện và đánh giá 4 mô hình (OLS cơ bản, OLS chọn lọc biến, Ridge, Lasso) trên tập dữ liệu thực.
    \end{itemize} \vspace{-0.2cm} & 
    \centering 100\% \tabularnewline \hline
    
\end{longtable}
